% Source: zcash/zips commit 2bdf39494320e5e04057ac93b1794031ebfb00b5, protocol/protocol.tex lines 8941-8957
% Spec version: v2026.7.0-344-g2bdf39 [NU6.2], section 5.2 "Bit layout diagrams"
% Pinned: 2026-10-24

\extralabel{boxnotation}{\lsubsection{Bit layout diagrams}{bitlayout}}

We sometimes use bit layout diagrams, in which each box of the diagram represents
a sequence of bits. Diagrams are read from left to right, with lines read from
top to bottom; the breaking of boxes across lines has no significance.
The bit length $\ell$ is given explicitly in each box, except when it is obvious
(e.g. for a single bit, or for the notation $\zeros{\ell}$ representing the sequence
of $\ell$ zero bits, or for the output of $\LEBStoOSP{\ell}$).

The entire diagram represents the sequence of \emph{bytes} formed by first
concatenating these \bitSequences, and then treating each subsequence of 8 bits
as a byte with the bits ordered from \emph{most significant} to
\emph{least significant}. Thus the \emph{most significant} bit in each byte
is toward the left of a diagram. \sapling{(This convention is used only in
descriptions of the \Sprout design; in the \SaplingAndOrchard additions,
\bitSequence/\byteSequence conversions are always specified explicitly.)} Where bit fields
are used, the text will clarify their position in each case.
